\subsection{An explicit change of first seed passes every field-norm level}

\begin{statement}{A.162}{Proved}
{A repaired two-level seed with compatible local norms and every finite global field-norm test}
Let \(d=649\), \(a=3d^{24}\), and
\(K_n=\Q(\rho_n)\), where \(\rho_n^{3^n}=a\) and
\(\rho_{n+1}^3=\rho_n\), with the positive real roots.
Retain \(K_1=\Q(\alpha)\), \(K_2=\Q(t)\),
\(\rho_1=d^8\alpha\), \(\rho_2=d^2t\) and the actual element
\(\beta_2\) of A.156.  Put
\[
 \begin{gathered}
 g_{11}=2+\alpha,\qquad g_{59}=2+2\alpha-\alpha^2,
 \qquad h_q=q/g_q^3\quad(q=11,59),\\
 f=h_{11}h_{59}^{2},\qquad
 b_1=(1-\alpha)f^3,\qquad b_2=\beta_2 f.
 \end{gathered}
\tag{A.162.1}
\]
These are actual \(S\)-units for \(S=\{2,3,11,59,\infty\}\), and
\[
       N_{K_2/K_1}b_2=b_1,\qquad N_{K_1/\Q}b_1=-2.
\tag{A.162.2}
\]
For every place \(v\) of \(K_2\), the element \(b_2\) belongs to a
compatible local norm family through all the algebras
\(K_m\otimes_{K_2}(K_2)_v\), \(m\geq2\).
For every finite \(m\geq2\), there is an actual
\(y_m\in K_m^\times\) with
\[
                         N_{K_m/K_2}y_m=b_2.
\tag{A.162.3}
\]
Consequently the fixed triple \((-2,b_1,b_2)\) extends to a compatible
actual global norm family through every prescribed finite height.
The proof retains the norm maps on the same fixed tower.
It also gives the exact quotient measuring the remaining restriction to
completed \(S\)-unit representatives.  Nonemptiness of an infinite global
norm family, or of the completed \(S\)-unit norm fibre at every height,
is not asserted here.
\end{statement}

\begin{proof}
For \(x=A+B\alpha+C\alpha^2\), the rational norm is
\[
                  Nx=A^3+3B^3+9C^3-9ABC.
\tag{A.162.4}
\]
It gives \(Ng_{11}=11\) and \(Ng_{59}=59\), hence \(Nh_q=1\).
Each \(g_q\) is integral of prime norm \(q\), so it generates the
old degree-one prime over \(q\).  The other old prime has degree two.
Thus, writing these primes as \(p_{q,0},p_{q,2}\),
\[
 \operatorname{div}_{K_1}(h_q)=p_{q,2}-2p_{q,0}.
\tag{A.162.5}
\]
Their extensions in \(K_2/K_1\) are single totally ramified cubics.
A.156 supplies \(N\beta_2=1-\alpha\) and its divisor over two.
Since the norm of the scalar \(f\in K_1^\times\) is \(f^3\),
A.162.2 follows exactly.  In \(K_2\), the additional divisor of \(b_2\)
is
\[
 3P_{11,2}-6P_{11,0}+6P_{59,2}-12P_{59,0}.
\tag{A.162.6}
\]
At two its divisor is still the unique degree-one prime, and at three
its valuation is zero.

\emph{The exact tame norm group at either changed prime.}
Fix \(q=11\) or \(59\), and first take the quadratic-residue
completion \(F=(K_2)_{P_{q,2}}\), with residue field \(k=\F_{q^2}\).
Put \(d'=d/q\).  The uniformizer and the unit part of the actual
radicand are
\[
 \pi=t^2/q,\qquad v_F(t)=2,\qquad
 \rho_2=\pi^8u_q,\qquad
 u_q=(d')^{-8}\alpha^{-5}\in\mathcal O_F^\times.
\tag{A.162.7}
\]
For \(e=3^{m-2}\), the extension \(F(\rho_m)/F\) is one totally
ramified field of degree \(e\), because \(v_F(\rho_2)=8\) is prime
to three.  Choose integers \(A_e,B_e\) with
\(8A_e+eB_e=1\).  Its uniformizer
\(\varpi_e=\rho_m^{A_e}\pi^{B_e}\) has the literal norm
\[
                          N\varpi_e=\pi u_q^{A_e}.
\tag{A.162.8}
\]
The residue of a norm of a unit is its \(e\)-th power, since the
residue degree is one.  Conversely every residue \(e\)-th power
lifts to a unit norm: use a scalar lift for that residue and then use
surjectivity of \(x\mapsto x^e\) on the principal units of \(F\).
The latter holds because \(e\) is prime to \(q\).
It follows, with no Galois assumption on this field extension, that
\[
 \pi^v c\in N(F(\rho_m)^\times)
 \quad\Longleftrightarrow\quad
 \overline c\,\overline u_q^{-A_ev}\in(k^\times)^e,
 \qquad c\in\mathcal O_F^\times.
\tag{A.162.9}
\]
Here \(v_3(q^2-1)=1\), so \((k^\times)^e=(k^\times)^3\) at every
positive depth, and \(A_e\equiv2\pmod3\).
This is the explicit comparison with the stable local norm test.

At its own quadratic place, \(h_q\) has valuation three and
\[
             h_q/\pi^3=(d')^{-4}\alpha^{-2}g_q^{-3}.
\tag{A.162.10}
\]
Rational units in \(\F_q^\times\) are cubes, because
\(q\equiv2\pmod3\).  Hence the cubic character of this normalized
unit is exactly that of \(\alpha\).
At the other changed prime \(h_q\) is a unit with trivial cubic
character, by its definition as a rational unit times a cube.
For complete signs, use the literal residue fields of A.156 and the
following values:
\[
\begin{array}{c|c|c|c|c|c}
q&\overline\beta_2&\omega_q=\lambda_q(\beta_2)
 &\lambda_q(\alpha)&v_F(b_2)&\overline{b_2/\pi^{v_F(b_2)}}\\\hline
11&6+3\alpha&10+6\alpha&\omega_{11}^{2}&3&8\\
59&23+57\alpha&58+6\alpha&\omega_{59}&6&29+6\alpha
\end{array}
\tag{A.162.11}
\]
The residue equations are \(\alpha^2-2\alpha+4=0\) at eleven and
\(\alpha^2-10\alpha+41=0\) at fifty-nine.
Both last-column entries have trivial cubic character.
Moreover both valuations are divisible by three, so the extra term
\(\overline u_q^{-A_ev}\) in A.162.9 also has trivial character.
Thus the fixed \(b_2\) is a norm at these places for every depth.
At the degree-one completion over \(q\), the same total-ramification
argument applies, now with residue field \(\F_q\), whose unit group
has order prime to three.  A.162.9 imposes no residue restriction there,
and proves that the target is a norm at every depth for its actual
valuation \(-6\) or \(-12\).

\emph{All other local places and local compatibility.}
At three \(F=(K_2)_3\) has odd degree nine over \(\Q_3\), contains
no \(\mu_3\), and satisfies \(v_F(\rho_2)=1\).
The exact local argument of A.140 applies to this pair \((F,\rho_2)\).
To retain its norm map, the splitting field of
\(X^{3^{m-2}}-\rho_2\) has group \(C_{3^{m-2}}\rtimes H\),
where \(H\) is the actual cyclotomic group and fixes the chosen root.
An element of \(H\) acts as \(-1\) modulo three; its commutator
kills the translation in the abelianization.  The image
\(H\to(C_{3^{m-2}}\rtimes H)^{\rm ab}\) is therefore onto.
In the local Artin norm diagram, the map on the right is this inclusion
modulo the commutator, while the map on the left is the field norm.
Given a target, lift its bottom Artin class through \(H\), and correct
by a norm from the splitting field.  Norm transitivity gives an actual
preimage in the pure field, as in A.140.8--9 and
\cite[III, Theorem 3.5]{MilneCFT}.
Since these pure fields are totally ramified, a preimage of our unit
\(b_2\) is a unit.

At two the distinguished completion of \(K_2\) is \(\Q_2\).
Its odd-power map on units is bijective, so the root tower has a
compatible degree-one factor there.  Put the target on that factor
and one on the other factors.  At the other two completions over two,
\(b_2\) is a unit.  Every higher root algebra is unramified: its
radicand is a unit and its defining derivative is a unit.  The unit
norm from each unramified field factor is onto, giving local norm
preimages there.  At every finite place outside \(2,3,11,59\), the
same unramified-unit argument applies.  At the real place use the
chosen real factor and put one on all complex factors.
At each complex place of \(K_2\), extend its chosen embedding to the
entire root tower, choosing the roots successively.  In
\(K_m\otimes_{K_2}\C\simeq\C^{[K_m:K_2]}\), put the image of
\(b_2\) in that distinguished factor and one in every other factor.
The norm is the product of the coordinates, and these selected embeddings
make the choices compatible under every transition norm.

These local choices can be made compatible at every height.
At each ramified field place, prescribe the valuation of a preimage
using the residue-degree-one norm formula; the resulting coset of the
unit group is compact.  At an unramified place use the compact product
of unit groups.  At the distinguished place over two retain its chosen
degree-one factor and use units on all other factors.  Norms are
continuous.  The norm fibre at every finite top height is nonempty by
the preceding computations, and pushing down from that height gives
all lower compatibility equations.  Thus the closed conditions in the
product of these compact sets have the finite intersection property.
This gives the claimed local inverse-limit family.  The real family
was explicit and is already compatible.

\emph{The actual full-decomposition Hasse norm principle.}
Let \(m\geq3\) and put
\[
 \Omega_m=\Q(\mu_{3^m},\rho_m),\quad U_m=(\Z/3^m)^\times,
 \quad G_m=(9\Z/3^m)\rtimes U_m,
 \quad H_m=U_m.
\tag{A.162.12}
\]
These are literally \(\operatorname{Gal}(\Omega_m/K_2)\) and
\(\operatorname{Gal}(\Omega_m/K_m)\).  Indeed \(v_3(a)=1\), so
\(X^{3^m}-a\) is Eisenstein over \(\Q_3\).
The non-Galois cubic field \(\Q_3(\sqrt[3]{a})\) cannot be contained
in an abelian cyclotomic extension.  Thus \(a\) is not a cube over
\(\Q_3(\mu_{3^m})\), and its Kummer order there is \(3^m\).
The full local degree is consequently \(3^m\varphi(3^m)\), which
also proves the global degree and the actual automorphisms fixing
\(\rho_m\) or multiplying it by a root of unity.
The translations fixing \(\rho_2=\rho_m^{3^{m-2}}\) are exactly
\(9\Z/3^m\), giving A.162.12.
In particular the decomposition subgroup over three is all of
\(G_m\).  Equivalently the local and global degrees agree; the
abelianization of the full holomorph is cyclotomic, so its maximal
unramified local quotient is trivial as well.

The norm-one torus and its exact character lattice are
\[
 T_m=\ker\bigl(R_{K_m/K_2}\mathbf G_m\xrightarrow{N}\mathbf G_m\bigr),
 \qquad X_m=\Z[G_m/H_m]/\Z\!\left(\sum_{gH_m}gH_m\right).
\tag{A.162.13}
\]
Hilbert 90 identifies its localization map with
\[
\begin{array}{ccc}
 H^1(K_2,T_m)&\simeq&K_2^\times/N_{K_m/K_2}K_m^\times\\
 \downarrow&&\downarrow\\
 \displaystyle\prod_vH^1((K_2)_v,T_m)&\simeq&
 \displaystyle\prod_v (K_2)_v^\times/
       N(K_m\otimes_{K_2}(K_2)_v)^\times.
\end{array}
\tag{A.162.14}
\]
The class of \(b_2\) lies in its kernel, by the established local
computations.  Torus duality identifies this kernel with the dual of
\[
 \ker\left(H^2(G_m,X_m)\longrightarrow
                      \prod_v H^2(D_{m,v},X_m)\right)
\tag{A.162.15}
\]
by \cite[Proposition 6.4]{BIR_Oki2026Norm}.
That proposition applies to every torus split by the displayed finite
Galois extension; no restriction to degree nine is used here.
The restriction at three in A.162.15 is the identity, so the kernel
is zero.

One can also see the precise inflation comparison directly.
A locally trivial lattice class restricts over \(\Omega_m\) to an
\((\Q/\Z)^{[K_m:K_2]-1}\)-valued character.  It is zero by local
triviality and Chebotarev.  Since
\(H^1(G_{\Omega_m},X_m)=0\), Hochschild--Serre then makes the
original class an inflation from \(H^2(G_m,X_m)\).
At the three-adic completion local inflation is injective by the same
\(H^1=0\) argument, and its decomposition group is \(G_m\).
This is exactly the map used in A.162.15.
It proves A.162.3.  For a prescribed finite height, choose its top
preimage and apply the actual norm maps downwards; A.162.2 attaches
\(b_1\) and \(-2\) at the bottom.
\end{proof}

\paragraph{The exact remaining support quotient.}
For \(i=2,m\), let \(E_i=\mathcal O_{K_i,S_i}^\times\), and let
\(P_i\) be the image of the divisor-outside-\(S_i\) map
\(K_i^\times\to\bigoplus_{\mathfrak p\notin S_i}\Z\mathfrak p\).
It is a free abelian group.  Thus
\(0\to E_i\to K_i^\times\to P_i\to0\) splits as an exact sequence
of abelian groups, and its pro-three completion is still exact.
The actual ideal norm makes the two rows a commutative diagram.
The snake map gives the canonical isomorphism
\[
 \ker\left(
 \frac{\widehat E_2^{(3)}}{N\widehat E_m^{(3)}}
 \longrightarrow
 \frac{\widehat{K_2^\times}^{(3)}}{N\widehat{K_m^\times}^{(3)}}\right)
 \simeq
 \frac{\ker(N:\widehat P_m^{(3)}\to\widehat P_2^{(3)})}
 {\operatorname{div}_{S^c}\ker
       (N:\widehat{K_m^\times}^{(3)}\to\widehat{K_2^\times}^{(3)})}.
\tag{A.162.16}
\]
Concretely choose a norm preimage of a target and take its divisor
outside \(S\); changing the preimage by a norm-one element gives
exactly the displayed denominator.  A.162.3 makes the completed class
of \(b_2\) lie in the left-hand kernel.  No assertion that its image
in the right-hand quotient vanishes has been proved.
The global field-norm fibres at every finite height are nonempty,
whereas the existence of a single compatible global family and the
vanishing of these fixed-support classes remain the explicit next
questions for this repaired seed.

\paragraph{Exact certificate.}
\path{code/two_inert_changed_seed_certificate.py}
verifies A.162.2 by rational polynomial arithmetic, the prime norms and
divisors, all residue values in A.162.11, and twenty-four explicit
Bezout/depth instances of A.162.8--9.  Its retained data contain the
complete coefficients of \(b_1,b_2\).  The all-depth assertions use
the displayed valuation, residue and compactness proofs; the finite
instances only check the literal formulas and signs.
